\documentclass{article}
\usepackage{amsmath}
\usepackage{amssymb}
\usepackage{amsthm}
\usepackage{mathtools}
\usepackage{geometry}
\usepackage{hyperref}

\newcommand{\rk}{\operatorname{rank}}
\newcommand{\GL}{\operatorname{GL}}
\newcommand{\R}{\mathbb{R}}
\newcommand{\Mat}{\mathrm{Mat}}

\theoremstyle{definition}
\newtheorem{theorem}{Theorem}[section]
\newtheorem{proposition}[theorem]{Proposition}
\newtheorem{lemma}[theorem]{Lemma}
\newtheorem{corollary}[theorem]{Corollary}
\newtheorem{definition}[theorem]{Definition}
\newtheorem{remark}[theorem]{Remark}

\title{Determinantal, Cyclic, and Critical Loci in Weight Space}
\date{}

\begin{document}
\maketitle

\section{Direct Weight-Space Varieties}

We now define the geometric objects needed for the saddle-combinatorics paper directly in the weight space $\Omega$, without passing through a residual parametrization.

\begin{definition}[Determinantal locus]
For $0\le r\le \min_\ell d_\ell$, define
\[
\Sigma_d^r
:=
\{W\in \Omega : \rk(\mu(W))=r\}.
\]
\end{definition}

\begin{definition}[Weight-space cyclic locus]
Fix a subset $S\subset \{1,\dots,r_{\max}\}$ with $|S|=r$, and let
\[
W_0=W_S:=\Sigma_{XY}P_Q.
\]
Define
\[
\Gamma_{d,S}
:=
\left\{
W\in \Omega
\;\middle|\;
\bigl(\nabla_{W_\ell}\mathcal L(W)\bigr)^\top=0
\text{ for every }1\le \ell\le L
\right\}
\]
\[
=
\left\{
W\in \Omega
\;\middle|\;
W_{<\ell}W_0W_{>\ell}=0
\text{ for every }1\le \ell\le L
\right\}.
\]
\end{definition}

\begin{definition}[Rank-$r$ cyclic locus]
The rank-$r$ part of the cyclic locus is
\[
\Gamma_{d,S}^r
:=
\Sigma_d^r\cap \Gamma_{d,S}.
\]
\end{definition}

\begin{definition}[Critical locus associated with $S$]
Define
\[
\mathcal C_{d,S}^r
:=
\left\{
W\in \Omega
\;\middle|\;
\mu(W)=P_SW^\ast
\text{ and }
\bigl(\nabla_{W_\ell}\mathcal L(W)\bigr)^\top=0
\text{ for all }1\le \ell\le L
\right\}.
\]
\end{definition}

\begin{proposition}[Equivalent description of the critical locus]
For fixed $r$ and $S$,
\[
\mathcal C_{d,S}^r
=
\left\{
W\in \Omega
\;\middle|\;
\mu(W)=P_SW^\ast
\text{ and }
W_{<\ell}W_0W_{>\ell}=0
\text{ for all }1\le \ell\le L
\right\}.
\]
In particular,
\[
\mathcal C_{d,S}^r \subseteq \Gamma_{d,S}^r.
\]
\end{proposition}

\begin{proof}
The equivalence follows from the gradient formula in the previous setup note and from the fact that, for critical points of rank $r$, the end-to-end map is equal to $P_SW^\ast$ for a unique subset $S$ of cardinality $r$.

Since $\mu(W)=P_SW^\ast$ has rank $r$, every point of $\mathcal C_{d,S}^r$ lies in $\Sigma_d^r$, and the cyclic equations place it in $\Gamma_{d,S}$. Hence $\mathcal C_{d,S}^r\subseteq \Gamma_{d,S}^r$.
\end{proof}

\begin{remark}[Why $\Gamma_{d,S}^r$ is larger than $\mathcal C_{d,S}^r$]
The variety $\Gamma_{d,S}^r$ only imposes:
\begin{itemize}
\item rank-$r$ on the end-to-end map;
\item the cyclic equations associated with $W_0=\Sigma_{XY}P_Q$.
\end{itemize}
The critical locus $\mathcal C_{d,S}^r$ imposes one more condition:
\[
\mu(W)=P_SW^\ast.
\]
Therefore the conjectural relation between $\Gamma_{d,S}^r$ and $\mathcal C_{d,S}^r$ cannot be literal equality. It must be expressed through an orbit correspondence under a group that is allowed to move the endpoint map.
\end{remark}

\begin{remark}[Comparison with the previous versions]
In \texttt{paper\_v0.tex} and \texttt{paper\_v1.tex}, the cyclic equations were encoded in a residual variety and then lifted to the critical locus through a block-normal-form slice. Here the order is reversed:
\begin{itemize}
\item first define the determinantal and cyclic loci directly in $\Omega$;
\item then formulate the relation with the critical locus directly in weight space.
\end{itemize}
This is the setup needed for the new saddle-combinatorics paper.
\end{remark}

\end{document}
